% Brief (tikz-diagram-craft), register row ch4-08.
% Reader question: through what states does one consent grant pass, and
%   what closes off its legitimate action set?
% Claim: a grant is ISSUED, then ACTIVE; each candidate action is ADMITTED
%   only if it clears a three-clause guard (in scope, under the stakes
%   ceiling, above the reversibility floor), else it ESCALATES -- both
%   outcomes return the grant to ACTIVE -- until TTL or revocation moves it
%   to EXPIRED/REVOKED, a terminal state that shrinks the daemons legitimate
%   action set.
% Evidence: Protocol Consent-grant lifecycle (Issue/Act/Witness/Expire),
%   Def. consent, whitepaper/legible-swarm.tex sec 4.2.2 (this pass).
% Must distinguish: the per-action guard on ADMIT (three clauses) from the
%   separate TTL/revoke guard that ends the grants life; that both ADMIT
%   and ESCALATE return to ACTIVE rather than ending the grant.
% Grammar chosen: state machine with a per-action guarded loop back to
%   ACTIVE, keyed to a ruled legend, reusing the swk workunit-machine spine.
% Grammar rejected: a flat four-step protocol ladder (issue, act, witness,
%   expire) would draw the four prose steps in order but hide that act is a
%   repeatable loop returning to ACTIVE, not a one-shot step between issue
%   and expiry.
% Five-second test: does a single escalated action end the grant, or does
%   the grant stay active for the next one?
\begin{figure}[H]
\centering
\pdfigurehue{pdteal}%
\begin{tikzpicture}[pd figure,x=1cm,y=1cm]
  \node[pd focus state,minimum width=2.4cm] (I) at (1.60,5.60) {issued};
  \node[pd focus state,minimum width=2.4cm] (A) at (1.60,3.60) {active};
  \node[pd terminal,minimum width=2.9cm]    (X) at (1.60,0.30) {expired / revoked};
  \node[pd ready state,minimum width=2.5cm] (D) at (7.90,4.80) {admitted};
  \node[pd warn state,minimum width=2.5cm]  (E) at (7.90,2.40) {escalated};
  \draw[pd focus arrow] (I) -- (A);
  \node[pd label,anchor=west] at (2.15,4.60) {issue};
  \draw[pd ready arrow] (A) to[bend left=12] (D);
  \draw[pd ready arrow] (D) to[bend left=12] (A);
  \draw[pd warn arrow] (A) to[bend right=12] (E);
  \draw[pd warn arrow] (E) to[bend right=12] (A);
  \draw[pd focus arrow] (A) -- (X);
  \node[pd badge] at (4.65,4.95) {1};
  \node[pd label,anchor=south] at (4.65,5.20) {action $x$: admit?};
  \node[pd label,anchor=north] at (4.65,2.75) {else};
  \node[pd badge] at (2.20,1.95) {2};
  \node[pd label,anchor=west,text width=2.4cm,align=left] at (2.50,1.95) {ttl or revoke};
  \def\ly{-1.55}
  \draw[pd rule] (0,\ly) -- (10.40,\ly);
  \node[pd title,anchor=north west] at (0.05,{\ly-.35}) {guard};
  \node[pd title,anchor=north west] at (5.20,{\ly-.35}) {what breaks without it};
  \draw[pd hairline] (0,{\ly-1.00}) -- (10.40,{\ly-1.00});
  \node[pd badge] at (.28,{\ly-1.45}) {1};
  \node[pd label,anchor=west,align=left,text width=4.5cm] at (.62,{\ly-1.45})
    {in scope, stakes$\le s_{\max}$,\\reversibility$\ge v_{\min}$};
  \node[pd label,anchor=west,align=left,text width=4.9cm,text=pderror] at (5.20,{\ly-1.45})
    {an out-of-scope or high-stakes\\action executes unsupervised};
  \node[pd badge] at (.28,{\ly-2.55}) {2};
  \node[pd label,anchor=west,align=left,text width=4.5cm] at (.62,{\ly-2.55})
    {ttl elapsed, or operator revokes};
  \node[pd label,anchor=west,align=left,text width=4.9cm,text=pderror] at (5.20,{\ly-2.55})
    {a stale or revoked grant keeps\\admitting actions};
\end{tikzpicture}
\caption{The consent-grant lifecycle of Def.~\ref{def:consent} and the
Consent-grant lifecycle protocol. A grant is issued, then active. Each
candidate action is admitted only if it clears the three-clause guard
(scope, stakes ceiling, reversibility floor); any clause it fails sends the
action to escalation instead. Both outcomes leave the grant active for the
next action -- a single escalated action does not end the grant. Only ttl
expiry or an operator revocation moves the grant to its terminal state,
after which the daemons legitimate action set shrinks by exactly that
grant. [internal]}
\label{fig:consent-lifecycle}
\end{figure}
